\documentclass[11pt]{article}
\usepackage{amsmath,amsthm,amssymb}
\usepackage[margin=1in]{geometry}
\usepackage{booktabs}
\usepackage{array}
\usepackage{stmaryrd}
\usepackage[colorlinks=true,linkcolor=blue,citecolor=blue,urlcolor=blue]{hyperref}

\newtheorem{theorem}{Theorem}
\newtheorem{proposition}[theorem]{Proposition}
\newtheorem{corollary}[theorem]{Corollary}
\newtheorem{lemma}[theorem]{Lemma}
\theoremstyle{definition}
\newtheorem{definition}[theorem]{Definition}
\newtheorem{remark}[theorem]{Remark}
\newtheorem{example}[theorem]{Example}

\newcommand{\Cat}{\mathsf{Cat}}
\newcommand{\Set}{\mathsf{Set}}
\newcommand{\Pos}{\mathbf{Pos}}
\newcommand{\Cont}{\mathsf{Cont}}
\newcommand{\Fam}{\mathrm{Fam}}
\newcommand{\Sub}{\mathrm{Sub}}
\newcommand{\cod}{\mathrm{cod}}
\newcommand{\nodes}{\mathrm{nodes}}
\newcommand{\branches}{\mathrm{br}}
\newcommand{\id}{\mathrm{id}}
\newcommand{\op}{\mathrm{op}}
\newcommand{\Tree}{\mathcal{T}}
\newcommand{\Path}{\mathcal{P}}
\newcommand{\Arb}{\mathcal{A}}
\newcommand{\Pfib}{\mathbb{P}}
\newcommand{\Uf}{\widehat{U}}
\newcommand{\inpr}{\hookrightarrow}
\newcommand{\Z}{\mathbb{Z}}
\newcommand{\Psip}{\Psi'}

\title{Branching is not a coproduct of branches:\\
\large unfolding versus descent in the container bifibration}
\author{MacBeth \\ \small Kodamai --- internal note (for Rick and Neil Ghani); CC Robin Langer}
\date{10 October 2026 (corrected revision)\\
\small corrects content commit \texttt{8972231}; content commit \texttt{4e7264d}}

\begin{document}
\maketitle

\begin{abstract}
The container bifibration $\Cont(\cod)=\Fam(\cod^\op)$ is a candidate categorical home for the
arboreal / game-comonad analysis of structure and power. In an earlier note we matched, over the
base of \emph{paths}, the arboreal relational fibration to a \emph{one-shape} container functor, up
to a fibrewise opposite. The obvious way to extend the match from paths to branching trees is to
make the branches of a tree into the \emph{shapes} of a container. We show this is wrong, and
precisely why: the branch assignment $\Psi(T)=(\branches(T),\{\mathrm{chain}_b\})$ computes the
colourings of the \emph{branch-unfolding} of the tree --- the disjoint union of its root-to-leaf
paths --- and so over-counts the colourings of the tree itself by the exact factor
$2^{\sum_v(\lambda(v)-1)}$, where $\lambda(v)$ is the number of branches through the node $v$. The
arboreal fibre reappears inside the container fibre as the sub-poset of \emph{prefix-consistent}
colourings: the object of descent data for the cover of the tree by its branches, an equalizer that
the shape-coproduct $\Fam=\Sigma$ cannot impose. The faithful encoding uses instead a
\emph{single} shape whose position set is the entire node set, with the tree order carried by the
\emph{base}; the encoding functor is \emph{contravariant}, $\Psip\colon\Tree^\op\to\Cont$, sending a
subtree embedding to its node injection read \emph{backward}, and the match is an isomorphism of
indexed posets $T\mapsto(2^{\nodes(T)},\subseteq)$. The one fact this needs --- that the arboreal
cartesian lift along a branching subtree embedding is computed by preimage --- is proved in six
lines inside our own category, with no arboreal axioms and no external citation. Branches do not
transform functorially at all: the object assignment $\Psi$ extends to no functor covering the
embeddings (a no-go, either variance); the relation between the branches of $T$ and of $T'$ is a
\emph{span}, which is exactly why gluing is descent and not a coproduct. The reusable principle:
shared sub-trajectories in a branching process are descent data, not a coproduct over branches.
\end{abstract}

\section{Introduction}\label{sec:intro}

\subsection*{The bridge, and the natural guess}

Two compositional vocabularies have grown up side by side. On one side, \emph{containers} and
\emph{directed containers} --- equivalently, small categories \cite{ahman-dcont-cat} --- package a
system as a set of \emph{shapes} together with, for each shape, a set of \emph{positions}; they
organise into the bifibration $\Cont(\cod)=\Fam(\cod^\op)$ whose fibres are predicates on positions
and whose two reindexing legs are direct and inverse image \cite{macbeth-contcod}. On the other
side, \emph{game comonads} and the \emph{arboreal} axiomatics \cite{arboreal-axioms} present
model-comparison games (pebble, Ehrenfeucht--Fra\"iss\'e, modal) as comonads whose coalgebras are
trees of plays, with a forgetful fibration $\Pfib\colon\Arb\to\Tree$ from coloured trees to trees.

That these two pictures should be \emph{the same picture} is a natural conjecture: both are theories
of ``a process together with a predicate on its states, pulled back along sub-trajectories''. In an
earlier note (Stage~1, \cite{macbeth-stage1}) we tested it at the smallest honest scale --- a single
unary colour, propositionally truncated, over the base $\Path$ of \emph{paths} (chains of plays) ---
and found a clean match, up to a twist. The arboreal path fibration $\Pfib|_\Path$ is isomorphic to
the \emph{fibrewise opposite} of a one-shape container functor, with arboreal relation-reflection
matching the \emph{opcartesian} (pullback) leg $\rho^*$ of $\Cont(\cod)$ --- not, as one might first
guess, the cartesian leg. (We will see that the honest way to say this is that the encoding functor
is \emph{contravariant}; \S\ref{sec:prelim}.)

Stage~1 reached only paths, because a one-shape container $(1,\{[n]\})$ \emph{is} a single chain of
positions and a branching tree is not. The registered next step --- Stage~2 --- conjectured the
obvious repair: to reach the full tree base $\Tree$, make the branches of a tree into the shapes of
a container,
\[
\Psi(T)=\bigl(\branches(T),\{\mathrm{chain}_b\}_{b\in\branches(T)}\bigr),
\]
one shape per root-to-leaf path, its positions the nodes along that path. The slogan was ``the
shape-coproduct $\Fam=\Sigma$ recovers the branch bookkeeping''.

\subsection*{The result}

The slogan is backwards, and the way in which it is backwards is the content of this note.

\begin{theorem}[Stage~2, the three-way dividing line]\label{thm:main}
In the unary, propositionally truncated regime, over the base $\Tree$ of finite rooted trees and
rooted-subtree embeddings:
\begin{enumerate}
\item[\textup{(1)}] \textup{(Positive, one shape --- proved, unconditional.)}
The arboreal relational fibration extends its Stage~1 match to \emph{all} of $\Tree$, branching
included, via the \emph{contravariant} one-shape functor $\Psip\colon\Tree^\op\to\Cont$,
$\Psip(T)=(1,\{\nodes(T)\})$, $\Psip(f)=(\id_1,j)$ with $j$ the node injection read backward: there
is an isomorphism of indexed posets $T\mapsto(2^{\nodes(T)},\subseteq)$, $f\mapsto j^{-1}$. The one
ingredient --- that the arboreal cartesian lift along a branching subtree embedding is preimage ---
is proved directly, in six lines, inside the coloured category itself (\S\ref{sec:positive}).
\item[\textup{(2)}] \textup{(Negative, branches-as-shapes --- proved.)}
The conjectured branch assignment $\Psi$ does \emph{not} give the match, at the level of
\emph{objects}: its container fibre is the Boolean lattice of colourings of the branch-unfolding
$\Uf(T)=\bigsqcup_b \mathrm{chain}_b$, strictly larger than the arboreal fibre by the factor
$2^{\sum_v(\lambda(v)-1)}$; the arboreal fibre reappears as the proper sub-poset of
\emph{prefix-consistent} colourings. Moreover $\Psi$ extends to \emph{no functor} covering the
embeddings in either variance (a no-go, \S\ref{sec:nogo}): branches do not transform functorially,
and the branch relation is a \emph{span}, not a map --- which is exactly why recovering the tree is
a descent step and not a coproduct.
\item[\textup{(3)}] \textup{(Arity --- the unary one-shape functor is blind to $k$-ary relations.)}
The specific unary functor $\Psip(T)=(1,\{\nodes(T)\})$ cannot see a $k$-ary relation for $k\ge2$,
and collapses distinct $k$-ary structures (on two nodes, all $16$ binary relations onto the same
unary datum, of which there are $4$). This is true and \emph{trivial}, and it is \emph{not} a
ceiling of the bridge: the $k$-th power container $\Psip_k(T)=(1,\{\nodes(T)^k\})$ reproduces the
$k$-ary fibre exactly (\S\ref{sec:arity}).
\end{enumerate}
\end{theorem}

The three clauses answer three different questions, and conflating them was the error in the
original conjecture. Clause~(2) is about the \emph{wrong} encoding of branching; clause~(1) about
the \emph{right} one; clause~(3) about an artefact of one particular functor, not an obstruction of
the bridge.

\subsection*{The mechanism: unfolding versus descent}

The heart is clause~(2), and it has a one-line diagnosis. The shape-coproduct $\Fam=\Sigma$ is a
\emph{coproduct}: applied to the branches of a tree it forms their disjoint union $\Uf(T)$, the
universal \emph{unfolding} of the tree by its paths, in which every shared prefix is duplicated once
per branch that passes through it. Colouring $\Uf(T)$ freely lets the duplicated copies of a node
disagree. The colourings of the \emph{tree} are exactly those that agree across copies --- the
\emph{prefix-consistent} ones --- and reconstructing the tree from its branch-cover by imposing that
agreement is a \emph{gluing}: a descent datum, the equalizer of the coproduct, of the wrong variance
for a coproduct to supply. Thus
\[
\Pfib\text{-fibre over }T
\;=\; \text{descent data for the branch-cover } \Uf(T)\twoheadrightarrow T,
\qquad
\Cont(\cod)\circ\Psi\text{-fibre}\;=\;\text{the raw cover colourings};
\]
they coincide iff the cover is trivial, i.e.\ $T$ is a path --- which is exactly why Stage~1 saw no
discrepancy. The honest encoding, clause~(1), sidesteps the unfolding altogether: keep \emph{one}
shape and put the whole node set in the positions, letting the \emph{base} $\Tree$ carry the tree
order through its subtree-embedding morphisms. One copy of each node, no duplication, nothing to
glue.

\subsection*{Why it matters}

The reusable statement is a modelling warning with teeth. Whenever a branching process --- a supply
chain forking into parallel routes, an agent-orchestration log splitting into concurrent
sub-dialogues, a provenance DAG --- is to be encoded in the container bifibration, the temptation is
to let each branch be a shape and take their coproduct. Our count shows precisely what that costs:
an over-representation by $2^{\sum_v(\lambda(v)-1)}$, one spurious bit of freedom for every extra
branch through every shared node. The shared history that distinguishes a tree from its unfolding is
\emph{descent data}, and the $\Fam=\Sigma$ coproduct structurally cannot see it. This is the
container-side form of our running theme that \emph{composition is a nerve-level-two datum}: the
sharing lives one level up from the branches themselves, and is recovered by an explicit gluing
step, never read off a product over branches.

\subsection*{Outline}

Section~\ref{sec:prelim} fixes the two fibrations and recalls the Stage~1 orientation in its
contravariant form. Section~\ref{sec:positive} gives the positive one-shape extension (clause~1) and
the six-line proof of its one ingredient, inside the coloured category, with no external citation.
Section~\ref{sec:unfolding} proves clause~(2)'s object-level over-count, with the V- and Y-tree
witnesses and the descent diagnosis. Section~\ref{sec:nogo} proves the no-go: $\Psi$ extends to no
functor. Section~\ref{sec:arity} treats clause~(3) and corrects the former ``ceiling''.
Section~\ref{sec:scope} records scope, open problems, and placement.

\paragraph{Note on this revision.}
This note corrects commit \texttt{8972231} following Rick's referee report of 10 October 2026. Three
things were wrong and are now fixed, all traceable to a single mistake --- the encoding functor was
made covariant and then patched by hand with a nearest-ancestor retraction $\rho_f$, which does not
compute preimage. We make $\Psip$ \emph{contravariant} and carry the node injection $j$ backward
(clause~1 becomes proved and unconditional); we replace the Jakl--Reggio dependence for the
cartesian-lift fact with a direct six-line proof (that paper is not needed, and is correctly cited
in \S\ref{sec:scope}); we state the no-go for $\Psi$ on morphisms explicitly; and we withdraw the
$k\ge2$ ``ceiling'', which is false.

\section{Two fibrations and the Stage~1 match}\label{sec:prelim}

Throughout we work in the \textbf{unary, propositionally truncated} regime: a single relation
symbol, interpreted as a \emph{subset} (a subobject, not a proof-relevant family). This is the
fragment in which the two theories have a chance of agreeing; \S\ref{sec:arity} isolates what the
unary one-shape functor cannot reach.

\subsection{The container bifibration}

Recall \cite{ahman-dcont-cat,gambino-kock} that a \emph{container} is a pair $(S,\{P_s\}_{s\in S})$
of a set $S$ of shapes and a family of position sets, and that containers organise into
$\Cont(\cod)=\Fam(\cod^\op)$, the bifibration over $\Cont(\Set)$ whose fibre over $(S,\{P_s\})$ is
$\prod_{s\in S}(\Set/P_s)^\op$. We use only its propositional truncation, recorded in
\cite{macbeth-contcod}. A container morphism $(S,\{P_s\})\to(S',\{P'_{s'}\})$ is a pair
$(\varphi\colon S\to S',\ \{\rho_s\colon P'_{\varphi s}\to P_s\})$: a \emph{forward} shape map and,
for each shape, a \emph{backward} position map. Then:

\begin{itemize}
\item \textbf{Fibre.} Over $(S,\{P_s\})$ the truncated fibre is
$\prod_{s\in S}\Sub(P_s)^\op=\bigl(2^{\bigsqcup_s P_s},\supseteq\bigr)$: predicates on the total
position set, ordered by \emph{reverse} inclusion.
\item \textbf{Cartesian leg.} The backward position map $\rho$ has cartesian reindexing
$(\Sigma_\rho)^\op$, truncated \emph{direct image} $\rho(-)\colon\Sub\to\Sub$.
\item \textbf{Opcartesian (pullback) leg.} The same $\rho$ has opcartesian reindexing $\rho^*$,
truncated \emph{inverse image} $\rho^{-1}$, left adjoint to the cartesian leg.
\end{itemize}

The two legs are the whole point: ``pull a predicate back along a sub-trajectory'' is $\rho^*$,
inverse image, \emph{not} the cartesian direction. Note that it is the \emph{opcartesian} leg that
performs inverse image; this will force the encoding functor below to be contravariant.

\subsection{The tree base and the arboreal relational fibration}

\begin{definition}\label{def:tree}
$\Tree$ is the category of finite rooted trees of bounded height (a tree is a finite poset with
least element $\mathrm{root}$ in which every down-set is a chain) and \emph{rooted-subtree
embeddings}: a morphism $f\colon T\to T'$ is an injection on nodes
$j=\nodes(f)\colon\nodes(T)\inpr\nodes(T')$ whose image is prefix-closed and contains the root. A
morphism of $\Tree$ is determined by its node map. $\Path\subset\Tree$ is the full subcategory of
chains (single branches); on $\Path$ the only morphisms are initial-segment inclusions --- exactly
the Stage~1 base.
\end{definition}

\begin{definition}\label{def:arb}
$\Arb$ is the category of trees carrying one unary colour $R\subseteq\nodes(T)$ --- the unary,
truncated relational lift of the arboreal coalgebra category, whose underlying coalgebras are the
copresheaves of \cite[Lemma~4.1]{macbeth-game}. A morphism $(T,R)\to(T',R')$ is a subtree embedding
$f\colon T\inpr T'$ that is \emph{colour-preserving}, $R\subseteq j^{-1}R'$ (equivalently
$j(R)\subseteq R'$), and there is at most one over each base map, so $\Arb\to\Tree$ is a poset
fibration. The functor $\Pfib\colon\Arb\to\Tree$ forgets the colour; it is faithful.
\end{definition}

Two facts about $\Pfib$ are all we use; both are elementary.

\begin{lemma}[arboreal fibre and reindexing]\label{lem:arbfibre}
\begin{enumerate}
\item A vertical map $R\to R'$ over $\id_T$ exists iff $R\subseteq R'$. Hence the $\Pfib$-fibre over
$T$ is the \emph{full} Boolean lattice $(2^{\nodes(T)},\subseteq)$ --- every subset is a legal
colouring, branching imposing no constraint.
\item Along a subtree embedding $f\colon T\inpr T'$ with node injection $j$, the arboreal reindexing
restricts the colour: $f^*R'=R'\cap\nodes(T)=j^{-1}(R')$ (relation-reflection). Inverse image is
strictly functorial.
\end{enumerate}
\end{lemma}

The identification in part~(1) depends on reading a node of $T$ as the unique play from the root to
it: the shared root is a \emph{single} play, so $\Pfib$ colours $\nodes(T)$ with the prefix shared,
not the branch-disjoint unfolding. This is the sentence the branch assignment $\Psi$ will violate.

\subsection{The Stage~1 orientation, contravariantly}

\begin{lemma}[Stage~1, \cite{macbeth-stage1}]\label{lem:stage1}
Present the Stage~1 encoding as a \emph{contravariant} functor $\Phi\colon\Path^\op\to\Cont(\Set)$,
$\Phi([n])=(1,\{[n]\})$, sending an initial-segment inclusion $[m]\inpr[n]$ to the container
morphism $(\id_1,j)\colon\Phi([n])\to\Phi([m])$ with backward position map the inclusion
$j\colon[m]\inpr[n]$. Then the $\Pfib$-fibre $(2^{[n]},\subseteq)$ and the truncated container fibre
$\Sub([n])^\op=(2^{[n]},\supseteq)$ are \emph{opposite} posets, and arboreal relation-reflection
$j^{-1}$ equals the \emph{opcartesian} leg $\rho^*$ of $\Cont(\cod)$ along $(\id_1,j)$ (differing
from the cartesian $(\Sigma_j)^\op=$ direct image). Hence
$\Pfib|_\Path\cong(\Phi^*\Cont(\cod))^{\mathrm{fib\text{-}op}}$ over $\Path^\op$.
\end{lemma}

We stress the variance, because it is where the earlier Stage~2 note went wrong. The backward map
that Stage~1 carries is the \emph{inclusion} $j$ (not a collapse): reindexing is $j^{-1}$, preimage.
The fibre-op and the $\rho^*$ identification used only that $j$ is an injection of sets and that
inverse image is functorial --- never that the positions formed a chain. This is why one of the two
candidate extensions survives verbatim, and the other does not.

\section{The positive extension: one shape, contravariant, positions the whole tree}%
\label{sec:positive}

\begin{definition}\label{def:psiprime}
$\Psip\colon\Tree^\op\to\Cont(\Set)$ sends $T$ to the \emph{one-shape} container
$(1,\{\nodes(T)\})$: a single shape, position set the entire node set as a \emph{flat} set, the tree
order discarded by $\Psip$ and carried by the base $\Tree$. On a subtree embedding
$f\colon T\inpr T'$ of $\Tree$, with node injection $j\colon\nodes(T)\inpr\nodes(T')$,
\[
\Psip(f)=(\id_1,\ j)\ \colon\ \Psip(T')\longrightarrow\Psip(T),
\]
where $j$ is the \emph{backward position map} of the container morphism (there is one shape, so a
morphism $\Psip(T')\to\Psip(T)$ is a single backward map $\nodes(T)\to\nodes(T')$, namely $j$). On a
chain, $\Psip([n])=(1,\{[n]\})=\Phi([n])$, so $\Psip$ extends $\Phi$; and $\Psip$ is functorial
because node injections compose functorially, $j_{g\circ f}=j_g\circ j_f$.
\end{definition}

\begin{remark}[Why the former $\rho_f$ was wrong]\label{rem:rhofwrong}
The earlier version made $\Psip$ covariant, $\Psip(f)=(\id_1,\rho_f)$ with
$\rho_f\colon\nodes(T')\twoheadrightarrow\nodes(T)$ the nearest-ancestor retraction. Neither
reindexing leg along $\rho_f$ equals relation-reflection $j^{-1}$. Concretely on the V-tree
$T'=\{r,a,b\}$, $T=\{r,a\}$ with $\rho_f(b)=r$: for $R'=\{b\}$, $j^{-1}R'=\varnothing$ but the direct
image $\Sigma_{\rho_f}R'=\{r\}$; for $R'=\{a,b\}$, $j^{-1}R'=\{a\}$ but $\Sigma_{\rho_f}R'=\{r,a\}$.
The opcartesian leg $\rho_f^{-1}$ even goes the wrong way (from $T$ to $T'$). Carrying $j$ backward,
as in Definition~\ref{def:psiprime}, is the fix, and it is the same orientation Stage~1 already
used.
\end{remark}

\subsection{The one ingredient, in six lines}

To promote a fibrewise match to an isomorphism of indexed posets we need that the arboreal
\emph{cartesian lifts} along subtree embeddings are computed by preimage. We prove this directly in
the coloured category $\Arb$, with no arboreal axioms and no external input.

\begin{lemma}[the branching cartesian lift is preimage]\label{lem:cartlift}
In $\Arb$ (Definition~\ref{def:arb}), let $f\colon(T,R)\to(T',R')$ be over the node injection $j$,
and recall $\Pfib$ is faithful. Then:
\begin{enumerate}
\item If $R=j^{-1}R'$ then $f$ is $\Pfib$-cartesian.
\item If $f$ is $\Pfib$-cartesian then $R=j^{-1}R'$.
\item $f$ is a pathwise embedding \textup{(}strong on each root chain\textup{)} iff $R=j^{-1}R'$.
\end{enumerate}
Consequently the cartesian lift of $f\colon T\inpr T'$ at $R'$ is $(T,j^{-1}R')$ --- preimage --- for
\emph{every} $T$, branching included.
\end{lemma}

\begin{proof}
(1) Given $g\colon(S,Q)\to(T',R')$ over $k=j\circ h$: then $j(h(Q))=k(Q)\subseteq R'$, so
$h(Q)\subseteq j^{-1}R'=R$; hence $h$ lifts to a morphism $(S,Q)\to(T,R)$, uniquely since $\Pfib$ is
faithful. (2) The map $(T,j^{-1}R')\to(T',R')$ over $j$ is a morphism; cartesianness of $f$ lifts
$\id_T$, giving $j^{-1}R'\subseteq R$. The reverse inclusion holds because $f$ is a morphism
($j(R)\subseteq R'$, i.e.\ $R\subseteq j^{-1}R'$). (3) ``Strong on each root chain $c$'' means
$R\cap c=j^{-1}R'\cap c$; the root chains cover $T$, so this is $R=j^{-1}R'$.
\end{proof}

The argument is uniform in $T$: there is no ``non-path case'', so no arboreal machinery is invoked.
(Exhaustive confirmation: all coloured trees with $\le4$ nodes, $3147$ morphisms, have
cartesian $\Leftrightarrow$ preimage $\Leftrightarrow$ pathwise in every case; our
\texttt{scratch/crown\_stage2\_repair\_verify.py} reproduces the V-tree slice.) Equivalently,
$\Pfib\cong\nodes^*(\Sub)$: the predicate fibration $\Sub\to\Set$ (fibre $(2^X,\subseteq)$,
reindexing $=$ preimage) pulled back along $\nodes\colon\Tree\to\Set$, $T\mapsto\nodes(T)$,
$f\mapsto j$.

\subsection{The match, as an isomorphism of indexed posets}

\begin{proposition}[one-shape extension]\label{prop:clause1}
In the unary truncated regime there is an isomorphism of indexed posets over $\Tree^\op$,
\[
\Pfib\;\cong\;\bigl(\Psip{}^*\Cont(\cod)\bigr)^{\mathrm{fib\text{-}op}}_{\mathrm{opfib}},
\qquad
T\longmapsto(2^{\nodes(T)},\subseteq),\quad f\longmapsto j^{-1},
\]
obtained from the \emph{opfibration part} of the pullback $\Psip{}^*\Cont(\cod)$ over $\Tree^\op$
after the fibrewise opposite. The branching of $T$ is encoded entirely in the base and its
subtree-embedding morphisms, while the container functor keeps one flat position set.
\end{proposition}

\begin{proof}
Fibrewise: over $\Psip(T)=(1,\{\nodes(T)\})$ the truncated container fibre is
$\Sub(\nodes(T))^\op=(2^{\nodes(T)},\supseteq)$, whose fibrewise opposite is the $\Pfib$-fibre
$(2^{\nodes(T)},\subseteq)$ of Lemma~\ref{lem:arbfibre}(1), on the nose --- one shape, one copy of
each node, no over-count. On morphisms: along $\Psip(f)=(\id_1,j)$ the opcartesian leg of
$\Cont(\cod)$ is $j^{-1}=$ inverse image, which equals arboreal relation-reflection
(Lemma~\ref{lem:arbfibre}(2)), and Lemma~\ref{lem:cartlift} identifies the arboreal cartesian lift
with the same preimage. The identification uses only that $j$ is a set injection and that inverse
image is functorial; no step uses a chain order. Composing gives the stated isomorphism.
\end{proof}

\begin{remark}[Well-typedness; nothing is conditional]\label{rem:welltyped}
The earlier phrase ``$\Pfib\cong(\Psip{}^*\Cont(\cod))_{\mathrm{fib\text{-}op}}$ as fibrations over
$\Tree$'' was ill-typed, because the variance is wrong: the encoding is contravariant, so the base
is $\Tree^\op$ and the correct object is an isomorphism of indexed posets $\Tree^\op\to\Pos$, as
above. With $\Psip$ contravariant, the word \emph{conditional} leaves the statement entirely:
clause~(1) is proved by Lemma~\ref{lem:cartlift} together with Proposition~\ref{prop:clause1}, with
no appeal to any external cartesian-lift theorem.
\end{remark}

\section{Branches as shapes unfolds the tree}\label{sec:unfolding}

This section proves the object-level half of clause~(2): the branch assignment computes the
colourings of the wrong object. Nothing here is conditional. (That $\Psi$ is moreover not a functor
is \S\ref{sec:nogo}.)

\begin{definition}\label{def:psi}
$\Psi$ sends a tree $T$ to the container $(\branches(T),\{\mathrm{chain}_b\}_b)$: one shape per
root-to-leaf branch $b$, with position set over $b$ the chain of nodes along $b$. On a chain there
is one branch and $\Psi$ agrees with $\Phi=\Psip$, so Stage~1 holds; off $\Path$ they diverge. We
treat $\Psi$ as an \emph{object assignment} here; \S\ref{sec:nogo} shows it is no functor.
\end{definition}

\begin{definition}\label{def:unfolding}
The \emph{branch-unfolding} of $T$ is $\Uf(T)=\bigsqcup_{b\in\branches(T)}\mathrm{chain}_b$, the
disjoint union of the branch-chains, with \emph{gluing map}
$\pi\colon\Uf(T)\twoheadrightarrow\nodes(T)$, $(b,v)\mapsto v$. For a node $v$, write $\lambda(v)$
for the number of branches through $v$ (equivalently, the number of leaves above $v$), so
$\pi^{-1}(v)$ has size $\lambda(v)$.
\end{definition}

\begin{lemma}[the exact over-count, an objectwise identity]\label{lem:overcount}
The truncated container fibre over $\Psi(T)$ is $\prod_b\Sub(\mathrm{chain}_b)^\op=(2^{\Uf(T)},
\supseteq)$, of cardinality $2^{|\Uf(T)|}$ with
\[
|\Uf(T)|=\sum_b|\mathrm{chain}_b|=\sum_{v\in\nodes(T)}\lambda(v),
\qquad
|\Uf(T)|-|\nodes(T)|=\sum_v(\lambda(v)-1).
\]
The pullback $\pi^*\colon 2^{\nodes(T)}\inpr 2^{\Uf(T)}$ is injective with image the
\emph{prefix-consistent} colourings --- those constant on each fibre $\pi^{-1}(v)$. Hence the
$\Pfib$-fibre $(2^{\nodes(T)},\subseteq)$ embeds as the proper sub-poset of prefix-consistent
colourings inside the $\Cont(\cod)\circ\Psi$-fibre $(2^{\Uf(T)},\supseteq)$, the excess being the
factor $2^{\sum_v(\lambda(v)-1)}$. The inclusion is proper exactly when some node has
$\lambda(v)\ge 2$, i.e.\ exactly when $T$ branches.
\end{lemma}

\begin{proof}
The fibre formula is the container fibre with $\bigsqcup_s P_s=\bigsqcup_b\mathrm{chain}_b=\Uf(T)$.
Counting positions by the node they glue to gives $|\Uf(T)|=\sum_v\lambda(v)$. A subset
$A\subseteq\Uf(T)$ equals $\pi^*(B)$ iff $A$ is a union of fibres of $\pi$ (constant on each
$\pi^{-1}(v)$); such $B=\{v:\pi^{-1}(v)\subseteq A\}$ is then unique, so $\pi^*$ is injective with
the stated image. The complementary degrees of freedom number $\prod_v2^{\lambda(v)-1}
=2^{\sum_v(\lambda(v)-1)}$, positive iff some $\lambda(v)\ge2$.
\end{proof}

This is a one-line count; it is a comparison of fibres \emph{over single objects}, and it holds
regardless of whether $\Psi$ acts on morphisms. In particular the two fibres have different
cardinalities as soon as $T$ branches, so there is \emph{no} fibrewise equivalence
$\Pfib\cong(\Psi^*\Cont(\cod))^\op$ --- not even onto the image of $\Psi$. (Verified on all $154$
labelled rooted trees with $\le6$ nodes; \texttt{scratch/crown\_stage2\_repair\_verify.py}.)

\begin{example}[the V-tree]\label{ex:v}
Let $V$ have root $r$ and two leaves $a,b$. Then $|\nodes|=3$, the branches are $ra,rb$, $|\Uf|=4$,
$\lambda(r)=2$, $\lambda(a)=\lambda(b)=1$, so $\sum_v(\lambda(v)-1)=1$. The arboreal fibre has
$2^3=8$ colourings; the $\Cont(\cod)\circ\Psi$-fibre has $2^4=16$. The $8$ excess colourings are the
\emph{inconsistent} ones --- colouring the root in branch $ra$ but not in $rb$, or vice versa.
\end{example}

\begin{example}[the Y-tree: deeper sharing]\label{ex:y}
Let $Y$ have root $r$, a single child $c$, and $c$ in turn with two children $x,y$. Then
$|\nodes|=4$, the branches are $rcx,rcy$, $|\Uf|=6$, $\lambda(r)=\lambda(c)=2$,
$\lambda(x)=\lambda(y)=1$, so $\sum_v(\lambda(v)-1)=2$: the whole shared spine $\{r,c\}$ is
double-counted. The arboreal fibre has $2^4=16$ colourings, the $\Cont(\cod)\circ\Psi$-fibre
$2^6=64$, ratio $2^2=4$.
\end{example}

\subsection{The diagnosis: $\Fam=\Sigma$ is unfolding, gluing is descent}

Lemma~\ref{lem:overcount} is a count; here is what it means. The branch family
$\{\mathrm{chain}_b\inpr T\}_b$ is a \emph{cover} of $T$ by its paths, for the \emph{Alexandrov
topology of down-sets}: each root chain is a down-set, hence open, and the chains cover $T$. Its
overlaps are $\mathrm{chain}_b\cap\mathrm{chain}_{b'}=\mathrm{chain}_{b\wedge b'}$, the shared prefix
down to the meet of $b$ and $b'$. A colouring of the tree is the same as a family of colourings of
the branches that \emph{agree on overlaps}: a \emph{descent datum} for the cover.

\begin{proposition}[descent description of the arboreal fibre]\label{prop:descent}
For every $T$, the image of $\pi^*$ --- the prefix-consistent colourings --- equals the set of
colourings of $\Uf(T)$ that agree on all pairwise overlaps of the branch cover; and this equals
$2^{\nodes(T)}$, the full arboreal fibre. Thus
\[
\Pfib\text{-fibre over }T
\;=\;\text{descent data (the \emph{equalizer}) for } \Uf(T)\twoheadrightarrow T,
\qquad
\Cont(\cod)\circ\Psi\text{-fibre}\;=\;\text{the raw cover colourings,}
\]
coinciding iff the cover is trivial, i.e.\ $T$ a path.
\end{proposition}

\begin{proof}
Two copies $(b,v)$, $(b',v)$ of a node $v$ lie in the overlap $\mathrm{chain}_{b\wedge b'}$ exactly
when both $b,b'$ pass through $v$; as $b,b'$ range over branches through $v$ these exhaust
$\pi^{-1}(v)$. Agreement on all overlaps therefore says a colouring is constant on each
$\pi^{-1}(v)$ --- prefix-consistency --- which by Lemma~\ref{lem:overcount} is the image of $\pi^*$,
in bijection with $2^{\nodes(T)}$. The gluing is the single equalizer
$2^{\nodes(T)}\rightarrowtail 2^{\Uf(T)}\rightrightarrows\prod_{b,b'}2^{\mathrm{chain}_{b\wedge b'}}$
(verified in \texttt{scratch/crown\_stage2\_repair\_verify.py}: consistent $=$ agree-on-overlap $=
2^{\nodes}$).
\end{proof}

The conceptual point is a mismatch of variance. $\Fam=\Sigma$ is a coproduct over shapes; fed the
branches it yields the disjoint union $\Uf(T)$, the unfolding of the tree. Reconstructing the tree
is the \emph{opposite} operation --- a gluing that identifies shared prefixes, the equalizer of the
coproduct, a limit-shaped datum. A coproduct cannot impose it. So $\Psi$ realises $\Pfib$ only as the
\emph{descent sub-fibration} of prefix-consistent colourings, never as the whole $\Cont(\cod)\circ
\Psi$.

\begin{remark}[reconciliation with clause~(1)]\label{rem:reconcile}
$\Psi$ and $\Psip$ differ exactly by unfolding versus gluing: $\Psi(T)$ is the branch cover,
$\Psip(T)$ is the tree. $\Cont(\cod)$ colours whatever flat position set it is handed. Hand it the
unfolded cover ($\Psi$) and it colours the cover, which descends to the tree only after an equalizer
$\Fam$ cannot supply; hand it the node set ($\Psip$) and it colours the tree directly.
\textbf{Branching is not a coproduct of branches}: the faithful encoding is a single shape whose
position set is the branching node set, with the base $\Tree$ carrying the order.
\end{remark}

\section{Branches do not transform functorially: a no-go}\label{sec:nogo}

Lemma~\ref{lem:overcount} compares fibres over single objects, and that is all the over-count needs.
But $\Psi$ was never defined on morphisms, and the question whether it \emph{can} be is not
cosmetic: it is the precise reason the branch picture is descent and not a functor. The answer is
no, in either variance.

\begin{theorem}[no functorial action for $\Psi$]\label{thm:nogo}
The object assignment $T\mapsto(\branches(T),\{\mathrm{chain}_b\})$ extends to \emph{no} functor
$\Tree\to\Cont$ nor $\Tree^\op\to\Cont$ whose action on morphisms covers the subtree embeddings.
Concretely:
\begin{enumerate}
\item \textup{(Contravariant fails on existence.)} A branch of the larger tree can restrict to a
non-branch prefix, leaving no valid shape image.
\item \textup{(Covariant fails on functoriality.)} A sibling-swap automorphism forces an impossible
fixed point of the branch shape-map.
\end{enumerate}
Equivalently, $\branches\colon\Tree\to\Set$ is not a functor compatible with the embeddings in
either variance.
\end{theorem}

Here ``covers $f$'' means the shape map of $\Psi(f)$ sends each branch to a branch compatible with
$j$: for the covariant form, a branch $b$ of the source with $j(\mathrm{chain}_b)\subseteq
\mathrm{chain}_{\varphi(b)}$; for the contravariant form, a branch $b'$ with
$j(\mathrm{chain}_{\varphi(b')})\subseteq\mathrm{chain}_{b'}$. Dropping compatibility makes $\Psi(f)$
carry no information about $f$, so it would not be ``the action of $\Psi$''; we return to that at the
end.

\begin{proof}
(1) On the V-tree inclusion $\{r,a\}\inpr\{r,a,b\}$: $\branches(T')=\{ra,rb\}$,
$\branches(T)=\{ra\}$. The shape map of a contravariant $\Psi(f)\colon\Psi(T')\to\Psi(T)$ would send
$b'=rb$ to a branch $b$ of $T$ with $j(\mathrm{chain}_b)\subseteq\mathrm{chain}_{rb}=\{r,b\}$; the
only candidate $\mathrm{chain}_{ra}=\{r,a\}$ maps to $\{r,a\}\not\subseteq\{r,b\}$. The branch $rb$
of $T'$ restricts to the prefix $\{r\}$, which is not a branch of $T$. No valid image.

(2) Covariant shape maps $\varphi\colon\branches(T)\to\branches(T')$ with
$j(\mathrm{chain}_b)\subseteq\mathrm{chain}_{\varphi(b)}$ exist per morphism (extend a branch to some
leaf above), but functoriality fails. Let $T=\{r,a\}$ ($a$ a leaf) and $T'=\{r,a,c,d\}$ with $c,d$
both children of $a$; let $f\colon T\inpr T'$ be the inclusion and $\sigma\colon T'\to T'$ the
automorphism swapping $c\leftrightarrow d$ (fixing $r,a$). Then $\branches(T)=\{\beta_a\}$ with
$\mathrm{chain}_{\beta_a}=\{r,a\}$, and $\branches(T')=\{\beta_c=rac,\ \beta_d=rad\}$.
Compatibility for $f$ forces $\varphi_f(\beta_a)\in\{\beta_c,\beta_d\}$ --- two choices.
Compatibility for $\sigma$ forces $\varphi_\sigma(\beta_c)=\beta_d$ and $\varphi_\sigma(\beta_d)=
\beta_c$, since $\sigma(\mathrm{chain}_{\beta_c})=\{r,a,d\}=\mathrm{chain}_{\beta_d}$: the swap, with
no fixed point. But $\sigma\circ f=f$ in $\Tree$ (equal node maps, as $\sigma$ fixes $r,a$ and
morphisms of $\Tree$ are determined by their node maps), so functoriality forces
$\varphi_\sigma\circ\varphi_f=\varphi_f$, i.e.\ $\varphi_f(\beta_a)$ is a fixed point of
$\varphi_\sigma$. The swap has none. Contradiction.
\end{proof}

Consequently clause~(2) of Theorem~\ref{thm:main} is honestly a statement about the \emph{object}
assignment $\Psi$. The genuine structure relating the branches of $T$ and of $T'$ is a
\emph{span} / correspondence
\[
R_f\subseteq\branches(T)\times\branches(T'),\qquad
(b,b')\in R_f\iff j(\mathrm{chain}_b)\subseteq\mathrm{chain}_{b'},
\]
left-total but neither left- nor right-functional --- which is precisely why $\Fam=\Sigma$ (a
functor) cannot realise it, and why recovering $\Pfib$ from $\Cont(\cod)\circ\Psi$ needs a
\emph{descent} step rather than a functorial map. The no-go is the container-side face of
\emph{composition is a degree-two datum}: sharing of branch-prefixes is a relation, not a map.

\begin{remark}[dropping compatibility does not rescue a functor]\label{rem:dropcompat}
If one lets the shape map ignore $f$, the only escape in~(2) is $\varphi_\sigma=\id$ --- but then
$\Psi(\sigma)$ does not cover the nontrivial automorphism $\sigma$ (it would identify the distinct
subtrees $r\to a\to c$ and $r\to a\to d$), so $\Psi$ is not an encoding of $\Tree$. Any functor that
distinguishes $\sigma$ from $\id$ on branches runs into the same contradiction.
\end{remark}

\section{The arity clause, and why there is no ceiling}\label{sec:arity}

Everything above is $k=1$. Clause~(3) concerns $k\ge2$, and here we must correct the earlier note,
which claimed a ceiling ``beyond $\Cont(\cod)$''. There is none.

\begin{proposition}[the unary one-shape functor is blind to $k$-ary relations]\label{prop:arity}
A $k$-ary relation on $T$ is a subset of $\nodes(T)^k$. For $k\ge2$ this is not a subobject in
$\Set/\nodes(T)$, and $\Cont(\cod)$ over $\Psip(T)=(1,\{\nodes(T)\})$ has fibre $\Sub(\nodes(T))^\op$
--- unary predicates only. So this particular functor ($2^n$ on $n$ nodes) collapses distinct
$k$-ary structures: on two nodes $\{0,1\}$ the relations $\varnothing$ and $\{(0,1)\}$, among all
$2^4=16$ binary relations, map to the same datum, of which there are only $2^2=4$. This statement is
true and \emph{trivial} --- a functor into unary predicates has no $k$-ary values.
\end{proposition}

\begin{proposition}[the $k$-th power container reproduces the $k$-ary fibre]\label{prop:power}
Let $\Psip_k\colon\Tree^\op\to\Cont$, $\Psip_k(T)=(1,\{\nodes(T)^k\})$, $\Psip_k(f)=(\id_1,j^k)$ with
$j^k=j\times\cdots\times j\colon\nodes(T)^k\to\nodes(T')^k$ the coordinatewise backward map. Then its
fibre over $T$ is $\Sub(\nodes(T)^k)^\op=(2^{\nodes(T)^k},\supseteq)$, and the opcartesian leg along
$\Psip_k(f)$ is $(j^k)^{-1}$ --- $k$-ary relation reflection,
$\{(v_1,\dots,v_k):(jv_1,\dots,jv_k)\in R'\}$. This is Lemma~\ref{lem:cartlift} verbatim with
$\nodes(T)$ replaced by $\nodes(T)^k$, so $\Psip_k$ reproduces the arboreal $k$-ary fibre exactly.
\end{proposition}

\begin{proof}
The fibre and the opcartesian leg are the \S\ref{sec:prelim} formulas with position set
$\nodes(T)^k$; the proof of Lemma~\ref{lem:cartlift} used only that positions form a set and that
inverse image is functorial. (Verified: $(j^2)^{-1}$ equals relation reflection on all sampled
binary relations on the V-tree, and the Cont fibre $2^{n^2}=16$ \emph{is} the arboreal $2$-ary
fibre; \texttt{scratch/crown\_stage2\_repair\_verify.py}.)
\end{proof}

So the former clause~(3) ``ceiling'' is \emph{false}, and we withdraw it. For a signature $\sigma$,
take one shape per relation symbol $r$ with positions $\nodes(T)^{\mathrm{ar}(r)}$; a coproduct over
\emph{symbols} is correct here, because distinct symbols are independent (this is a genuine
coproduct, unlike the spurious branch-coproduct of \S\ref{sec:unfolding}). If the intended $k$-ary
fibre is the Ehrenfeucht--Fra\"iss\'e one --- related tuples must be pairwise comparable, on a common
root chain --- take the positions to be the \emph{comparable} $k$-tuples; then the earlier count
$2^{n^k}$ is also wrong and should read $2^{\#\{\text{comparable }k\text{-tuples}\}}$. In every case
the $k$-ary structure sits \emph{inside} $\Cont(\cod)$, at the container with positions
$\nodes(T)^k$; what survives of clause~(3) is only the trivial observation that the \emph{unary}
$\Psip$ cannot see it.

\section{Scope, open problems, and placement}\label{sec:scope}

\paragraph{What is proved.}
Clause~(1) --- the positive one-shape extension over $\Tree$, via the \emph{contravariant}
$\Psip\colon\Tree^\op\to\Cont$ with $j$ carried backward --- is \emph{proved, unconditional}
(Proposition~\ref{prop:clause1}), its one ingredient discharged by the six-line
Lemma~\ref{lem:cartlift} with no external citation. Clause~(2)'s over-count $2^{\Uf}\supsetneq
2^{\nodes}$ with factor $2^{\sum_v(\lambda(v)-1)}$, and its descent reading, are proved
(Lemma~\ref{lem:overcount}, Proposition~\ref{prop:descent}, Examples~\ref{ex:v},~\ref{ex:y}), as a
statement about the \emph{object} assignment $\Psi$; and $\Psi$ extends to no functor
(Theorem~\ref{thm:nogo}). Of clause~(3), only the trivial statement survives
(Proposition~\ref{prop:arity}); the ceiling is withdrawn (Proposition~\ref{prop:power}). The
objectwise over-count has a machine-checked Lean oracle (the registry node
\texttt{lean-crown-stage2-overcount-identity}).

\paragraph{On Jakl--Reggio.}
The earlier version attributed the cartesian-lift fact to T.~Jakl and L.~Reggio, \emph{On the Axioms
of Arboreal Categories}, arXiv:2603.21841 (CMCS 2026), Theorems~23 and~26. We do \emph{not} use it:
the fact has the direct six-line proof of Lemma~\ref{lem:cartlift} inside our own coloured category
$\Arb$, needing no arboreal axioms (which, moreover, we never verified for $\Arb$ --- our base
$\Tree$ is rooted, exactly the case where the older connectedness axiom fails). The citation is
corrected here and recorded as not needed.

\paragraph{Open problems.}
\begin{enumerate}
\item \textbf{Proof-relevance.} The match is propositional. The full container fibre
$(\Set/\nodes)^\op$ is strictly richer than $\Sub(\nodes)^\op$; a proof-relevant comparison would
need an arboreal theory with witness-set colours.
\item \textbf{The descent completion.} Proposition~\ref{prop:descent} is elementary (a single
equalizer per tree). A statement ``$\Pfib$ is the descent completion of $\Cont(\cod)\circ\Psi$'' can
\emph{not} be phrased as a comparison of functors, since $\Psi$ is not one
(Theorem~\ref{thm:nogo}); it must be stated as descent along the \emph{span} $R_f$. That is the
honest form of the open problem.
\item \textbf{Beyond embeddings.} $\Tree$ here has subtree embeddings only; open maps and non-monic
tree morphisms are untreated, inherited from the Stage~1 conventions.
\end{enumerate}

\paragraph{Placement.}
The net correction to the registered Stage~2 conjecture is clean. It is \emph{not} ``$\Psi$ positive
on the unary fragment, failing at $k\ge2$''. It is that the branch-multishape $\Psi$ fails already
at $k=1$ \emph{branching} (unfolding versus descent, and not even a functor); the positive Stage~2
holds over the full tree base via the contravariant one-shape $\Psip$; and $k\ge2$ is no obstruction
at all, merely invisible to the \emph{unary} functor. The reusable principle, worth carrying into the
applied side of the programme: \emph{when encoding a branching process in the container bifibration,
do not make the branches into shapes.} That unfolds the process and over-counts shared history by
one bit per extra branch through each shared node; keep one shape and put the whole state set in the
positions, because the sharing --- shared provenance, shared sub-trajectory, composition itself ---
is a descent datum that the $\Fam=\Sigma$ coproduct structurally cannot see. This is the
container-side statement of ``composition is a degree-two datum'', and it applies verbatim to
supply-chain provenance models and to agent-orchestration logs: shared sub-trajectories must be
glued by an explicit descent step, not recovered from a product over branches.

\begin{thebibliography}{9}

\bibitem{ahman-dcont-cat}
D.~Ahman, J.~Chapman, T.~Uustalu,
\emph{Directed containers as categories},
in \emph{Proc.\ MSFP}, EPTCS, 2016.

\bibitem{gambino-kock}
N.~Gambino, J.~Kock,
\emph{Polynomial functors and polynomial monads},
Math.\ Proc.\ Cambridge Philos.\ Soc.\ \textbf{154} (2013), 153--192.

\bibitem{arboreal-axioms}
S.~Abramsky, L.~Reggio,
\emph{Arboreal categories: an axiomatic theory of resources},
Log.\ Methods Comput.\ Sci.\ \textbf{19} (2023).

\bibitem{jakl-reggio}
T.~Jakl, L.~Reggio,
\emph{On the axioms of arboreal categories},
arXiv:2603.21841 (CMCS 2026).
(Theorem~23: a morphism is Cartesian iff a pathwise embedding; Theorem~26: the path functor is a
Street fibration. \emph{Not used in this note}; the cartesian-lift fact is proved directly in
Lemma~\ref{lem:cartlift}.)

\bibitem{macbeth-contcod}
MacBeth,
\emph{The logic of containers: $\Cont(\cod)=\Fam(\cod^\op)$ as a bifibration},
Kodamai internal proof note, 28 Aug 2026.

\bibitem{macbeth-game}
MacBeth,
\emph{Game comonads are polynomial comonads} (Lemma~4.1: coalgebras of a forest comonad are
copresheaves),
Kodamai internal proof note, 1 Oct 2026.

\bibitem{macbeth-stage1}
MacBeth,
\emph{Refined Bridge-3 crown, Stage 1: the orientation gap --- relation-reflection is $\rho^*$,
not $(\Sigma_\rho)^\op$},
Kodamai internal proof note, 2 Oct 2026.

\end{thebibliography}

\end{document}
